\documentclass[11pt]{article}
\usepackage[margin=1in]{geometry}
\usepackage{amsmath,amssymb,amsthm,hyperref}
\newtheorem{definition}{Definition}
\newtheorem{lemma}{Lemma}
\newtheorem{obligation}{Open obligation}
\newcommand{\F}{\mathbb F}
\title{LSPV Candidate Spine: Canonical Boolean Separators and Reconstruction}
\author{LOWER\_BOUND-0001, P versus NP project}
\date{10 October 2026}
\begin{document}
\maketitle
\paragraph{Status and scope.}
This document isolates a constructive candidate from recovered LSPV work.
The algebraic density lemma below is proved; the canonical separator
theorem, seed-indexed reconstruction interface, and final diagonal remain
OPEN. No proof of $P\ne NP$ is claimed. The current GitHub ARC/WRAPPER
spine is a distinct route, treated in the companion manuscript.

\section{The exact object that requires compression}
Project sources include Drive MAIN LSPV-5::58501,
\url{https://docs.google.com/document/d/1z4Jaoi4P0ZJXU8DTemXW6Eqb3vvt5RosAZTWuVMeCG0/edit},
and LSPV-5::26501,
\url{https://docs.google.com/document/d/1MGtSvJDCXCOpA1ulYycq_S-iExLaSaUq7_ZeHn-8e3o/edit}.
The latter was reviewed through search hydration in this run; the former
was fetched in full. Current primary-source numbering is van Melkebeek
and Mocelin Sdroievski \cite{vms2025}, Section 4.1, Definition 4.1,
Lemma 4.2, and the commit-and-evaluate reconstruction in 4.4.
Earlier project documents use numbering from an older manuscript.

\begin{definition}[Source-fixed successful range]
Fix a source instance $x$, a source time scale $T=T(|x|)$,
one fully specified generator $H$ with its parameters, and one successful
nondeterministic path $\pi$ of $H(x)$. Let
\[
 R_{x,\pi}\subseteq\{0,1\}^{m}
\]
be the entire set of output strings on that path.
The path fixes one consistent input encoding, all chosen accepted proof
objects, and all their low-degree extensions.
\end{definition}
For the published PCP/PCPP construction, the relevant output is a union
of RMV outputs for the encoded input and the chosen proof for each output
bit. RMV itself includes recursively generated lines and all local seeds.
Separation of one local extractor slice, one restriction, or one proof
polynomial is not automatically separation of this full union.

\begin{definition}[Canonical Boolean separator target, Route A]
Fix $0<\delta<1/2$ and a charged size bound $B(T)$.
For every source instance required by the eventual endpoint, seek a
successful path $\pi$ and one Boolean circuit $D_{x,\pi}$ such that
\[
 \forall z\in R_{x,\pi}\quad D_{x,\pi}(z)=0,\qquad
 \Pr_{z\sim U_m}[D_{x,\pi}(z)=1]\ge \tfrac12+\delta,\qquad
 \operatorname{size}(D_{x,\pi})\le B(T).
\]
\end{definition}
The quantifier is $\forall x\ \exists\pi\ \exists D\ \forall z\in R_{x,\pi}$,
not $\forall z\ \exists D_z$. The downstream soundness quantifier
must still control every certificate the verifier permits.

The circuit may be guessed existentially if the complete endpoint
requires only a robust truth value. Its identity need not be recovered
deterministically unless a later stage consumes it. Encoded length,
verification, evaluation, and density checks remain charged.

\section{One proved representation lemma}
Let $\mathcal Q_m$ be the $\F_2$ vector space of reduced multilinear
polynomials of degree at most two in $z_1,\ldots,z_m$. Its coordinates
are $1,z_i,z_iz_j$ with $i<j$; its dimension is
$1+m+\binom m2$. For $R\subseteq\{0,1\}^m$, define
\[
 K_2(R)=\{Q\in\mathcal Q_m:\ \forall z\in R,\ Q(z)=0\},
 \qquad
 \mu(Q)=\mathbb E_{z\sim U_m}(-1)^{Q(z)}.
\]
For a two-dimensional subspace $L\le K_2(R)$, let $w(L)$ be the number
of monomial coordinates used by at least one member of $L$.
Define
\[
 J_{2,\delta}(R)=
 \min_{\substack{L\le K_2(R),\ \dim L=2\\
                 \sum_{0\ne Q\in L}\mu(Q)\le1-4\delta}}w(L),
\]
with value $+\infty$ when the feasible set is empty.

\begin{lemma}[Quadratic-OR separator]
If $J_{2,\delta}(R)=w<\infty$, then a Boolean separator $D$ rejects $R$,
accepts at least $1/2+\delta$ of the cube, and has bounded-fanin circuit
size $O(w+1)$. A direct sparse description uses
$O((w+1)\log(m+2))$ bits.
\end{lemma}
\begin{proof}
Choose a minimizing plane $L=\langle Q_1,Q_2\rangle$ and put
$D(z)=Q_1(z)\lor Q_2(z)$. Both quadratics vanish on $R$.
For every $z$,
\[
 {\bf1}[Q_1(z)=Q_2(z)=0]
 =\tfrac14(1+(-1)^{Q_1(z)})(1+(-1)^{Q_2(z)}).
\]
Averaging gives
\[
 \Pr[D=1]
 =\tfrac34-\tfrac14\bigl(\mu(Q_1)+\mu(Q_2)+\mu(Q_1+Q_2)\bigr)
 \ge\tfrac12+\delta.
\]
At most $w$ distinct monomials are needed. Each has at most two
variables and is computed with at most one AND gate. Each quadratic
is the XOR of a subset of these monomials. Bounded-fanin XOR and the
final OR need $O(w+1)$ gates over any usual complete Boolean basis.
Each monomial has at most two variable indices, each using
$O(\log(m+2))$ bits, plus constant bookkeeping.
\end{proof}
Circuit size and descriptor length are distinct.
In particular, $O(m)$ gates do not by themselves give a literal $O(m)$-bit
sparse coefficient description. For $m=T^\eta$, the logarithm may be
absorbed in $T^{\eta+o(1)}$, but only after it is charged.

\begin{obligation}[LSPV-A: canonical sparse, biased kernel]
For the exact $R_{x,\pi}$, prove the existence of a successful $\pi$
with $J_{2,\delta}(R_{x,\pi})$ within the charged support budget,
or construct another circuit family satisfying Route A.
\end{obligation}
This is sufficient for one representation class, not necessary for all
separators. Kernel dimension, low degree, or support size alone does
not imply the three-bias condition. The full union must be respected.

\section{An alternative graph target and its cost}
Suppose the complete range is indexed by a fully specified total map
\[
 G_{x,\pi}:\{0,1\}^{s}\longrightarrow\{0,1\}^{m}.
\]
All component indices, recursive-line choices, seeds, invalid-index
conventions, and their distribution must be defined in the actual
generator. A padding convention must not silently change density.

\begin{definition}[Route B graph separator]
Seek one circuit $D^*_{x,\pi}(y,z)$ satisfying
\[
 \forall y\quad D^*_{x,\pi}(y,G_{x,\pi}(y))=0,\qquad
 \Pr_{(y,z)\sim U_s\times U_m}[D^*_{x,\pi}(y,z)=1]\ge p_0,\qquad
 \operatorname{size}(D^*_{x,\pi})\le B^*(T),
\]
where $p_0>0$ is fixed and fits the error gap of a subsequently proved
reconstruction theorem.
\end{definition}
This differs from a seed-independent $D(z)$. The published
local-extractor definition does not, merely from its statement, give
the strong joint-seed reconstruction required here.

\begin{lemma}[Trivial graph separator and its limitation]
If $G$ has a circuit of size $S$, then
$D^*(y,z)={\bf1}[z\ne G(y)]$ has size $O(S+m)$ and joint density
$1-2^{-m}$.
\end{lemma}
\begin{proof}
Compute $G(y)$, XOR it coordinatewise with $z$, and OR the differences.
For every $y$, exactly one $z$ is rejected.
\end{proof}
This does not make $S$ small. Hard-wiring a $T$-scale proof polynomial
or recomputing it can consume the resource the construction intends
to save.

\begin{obligation}[LSPV-B: source-specific graph compression]
Prove a graph separator substantially cheaper than evaluating the full
source-specific graph, with its description, evaluation, and verifier
within the same declared budget.
\end{obligation}

\section{Corrected interpolation guard}
\begin{lemma}[Finite-point evaluation surjectivity]
For distinct points $P_1,\ldots,P_m\in\F_q^d$ and arbitrary
$a_1,\ldots,a_m\in\F_q$, there is a polynomial $p$ of total degree
at most $m-1$ satisfying $p(P_i)=a_i$ for all $i$.
\end{lemma}
\begin{proof}
For each ordered pair $i\ne j$, choose a coordinate $k(i,j)$ where
$P_i$ and $P_j$ differ and set
\[
 L_i(X)=
 \prod_{j\ne i}
 \frac{X_{k(i,j)}-(P_j)_{k(i,j)}}
      {(P_i)_{k(i,j)}-(P_j)_{k(i,j)}}.
\]
Each denominator is nonzero. Then $L_i(P_i)=1$ and
$L_i(P_j)=0$ for $j\ne i$. The sum $p=\sum_i a_iL_i$ works.
\end{proof}
Arbitrary degree-bounded source polynomials can therefore realize a
full output cube on sufficiently short distinct evaluation sequences.
No nonzero polynomial function on that cube vanishes everywhere.
The corresponding assertion about formal polynomials is false:
$X_1^q-X_1$ is nonzero formally but vanishes on all of $\F_q^m$.
Reduced representatives of individual degree below $q$ avoid the
ambiguity. This corrects wording in MAIN::58501. The coordinate-product
proof also avoids an unnecessary shared separating-functional assumption.

Source-specific structure is indispensable. The lemma concerns the
unrestricted polynomial class, not a proof that the canonical
accepted-proof class has a full range.

\section{Exact missing downstream obligations}
\begin{description}
\item[] \textbf{LSPV-SOURCE.}
Give the canonical source-to-PCPP-to-low-degree-extension map, generator
parameters, and full successful output union. Every object used in
completeness must refer to the same realization.
\item[] \textbf{LSPV-COMPRESSION.}
Prove Route A or B on that source family for every instance used by the
final diagonal. A toy low-degree family is not a replacement.
\item[] \textbf{LSPV-BINARY.}
If starting from an $\F_q$ orbit recurrence, derive a relation after
the exact binary code-coordinate map. An odd-characteristic field
recurrence is not automatically a relation on the exposed Boolean bits.
\item[] \textbf{LSPV-INTERFACE.}
Match Route A's circuit size, input length and density to the precise
published reconstruction parameters, including constants and error.
For Route B prove a modified joint-seed reconstruction theorem.
Per-seed distinguishers do not supply the required single circuit.
\item[] \textbf{LSPV-SOUNDNESS.}
If $D$ is guessed, prove completeness for at least one certificate
and false-endpoint soundness for every allowed certificate.
Charge density checking, oracles and randomness. Prove whether the
endpoint consumes only a bit or later needs $D$'s identity.
\item[] \textbf{LSPV-RESOURCES.}
Fix the hypothetical SAT decider and exponent $k$ once.
Account for circuit encoding, source access, reconstruction, queries,
derandomization, PH-collapse exponents, output and census.
The archive's schematic $k^4$ multiplier is not promoted here to a
proved exact compiler theorem.
\item[] \textbf{LSPV-DIAGONAL.}
State the exact source language/function, machine model, constructible
time bound, input/program length relation, uniformity, simulation and
self-reference overhead. Prove a strictly faster single algorithm
contradicting a named applicable hierarchy/diagonal theorem.
A sub-$T$ estimate without these data is not a contradiction.
\item[] \textbf{LSPV-BARRIERS.}
Identify the step that fails under legitimate common-oracle replacement.
Same-realization binding and Boolean syntax alone do not certify
nonrelativization or nonalgebrization.
\end{description}

\paragraph{Next falsification task.}
For one exact canonical successful path, construct its full RMV range
at feasible small parameters. Compute $K_2$, candidate minimum supports
and the three biases. Compare with accepted-proof padding and relabeling
controls. A finite counterexample can retire an overbroad claim.
A finite success remains computational evidence until a uniform
structural theorem is proved.

\paragraph{Literature coverage.}
The current journal PDF of \cite{vms2025} was read at the local-extractor,
RMV, and reconstruction interfaces. This document does not rederive its
complete reconstruction proof. The older author-PDF link returned 404;
the current author-hosted journal PDF works:
\url{https://pages.cs.wisc.edu/~dieter/Papers/r-am-instance-wise-cc.pdf}.
The Shaltiel--Umans 2005 source was located through Caltech, but its
separate PDF fetch timed out; no unverified transfer from it is used
as a proved premise here.

\bibliographystyle{plain}
\bibliography{lspv_candidate_refs}
\end{document}

